The pseudo-genetic formulation reframes evaluation from visibility-only indicators toward compositional analysis of knowledge flow. By separating inherited semantic mass from residual contribution, the framework provides interpretable evidence at both citation-edge and paper levels.

\subsection{How to read contribution and inheritance}
A high contribution factor indicates lower explainability by cited parents under the selected representation; it is not a direct claim of quality or impact. A high inheritance factor indicates stronger semantic dependence on predecessors; it is not a deficit signal. Synthesis papers, benchmarks, and consolidation work can have high scientific value despite high inherited mass.

\subsection{Analytical uses}
The framework supports:
\begin{itemize}
    \item lineage tracing for hypothesis-driven analysis,
    \item comparison of inheritance concentration across branches,
    \item detection of potential cross-domain transfer channels,
    \item and structured review of novelty claims against citation-supported dependence.
\end{itemize}

\subsection{Interpretive boundaries}
Outputs should not be used as deterministic judgments of researcher merit and should not be treated as causal evidence of intellectual influence without corroborating analysis. Strong claims should rely on convergence between quantitative outputs, uncertainty summaries, and expert qualitative inspection.
